\فصل{مقدمه}

در این فصل، مسئلهٔ پژوهش و اهمیت آن را شرح می‌دهیم، مروری کوتاه بر پیشینهٔ موضوع داریم، اهداف و پرسش‌های پژوهش را بیان می‌کنیم و در پایان، ساختار پایان‌نامه را معرفی می‌کنیم.


\قسمت{تعریف مسئله}

بازشناسی خودکار گفتار\پاورقی{Automatic Speech Recognition (ASR)} فرایند تبدیل سیگنال گفتار به متن است. دقت آن علاوه بر زبان و محتوا، به مقدار و تنوع دادهٔ آموزشی و شباهت شرایط آموزش به محیط استفاده بستگی دارد. در تماس‌های تلفنی و انتقال صدا بر بستر پروتکل اینترنت\پاورقی{Voice over IP (VoIP)}، سیگنال ممکن است هم‌زمان در معرض نویز، تغییر بهره، محدودیت پهنای‌باند، فشرده‌سازی کدک و افت بسته قرار گیرد؛ بنابراین، مدل آموزش‌دیده با گفتار نسبتاً پاک لزوماً در تماس واقعی همان دقت را ندارد.

این ناهمخوانی دامنه در فارسی با کمبود دادهٔ برچسب‌خوردهٔ متنوع همراه است. منابع فارسی محدودترند و بخش بزرگی از آن‌ها گفتار خوانده‌شده یا قطعه‌های کوتاه‌اند؛ حال‌آنکه گفتار واقعی می‌تواند خودانگیخته، پیوسته، نویزی یا باریک‌باند باشد. با وجود پیشرفت از فارس‌دَت و نویسا تا مدل‌های عصبی، پوشش سبک‌ها و کانال‌ها و ارزیابی خارج از دامنه همچنان چالش‌اند \مرجع{bijankhan1994farsdat,sameti2011large,veisi2020persian}. در آزمایش‌های نخست، نزدیک به هزار ساعت گفتار فارسی در قالبی مشترک آماده و نمونه‌های بلندتر نیز ساخته شد.

ویسپر\پاورقی{Whisper} به‌دلیل پیش‌آموزش با نظارت ضعیف روی داده‌های چندزبانه، نقطهٔ آغاز مناسبی برای انتقال به فارسی است \مرجع{radford2023robust}؛ اما پیش‌آموزش جای مواجههٔ مستقیم با تنوع هدف را نمی‌گیرد. مسئله آن است که در محدودیت داده و محاسبات، افزایش تنوع و مقیاس داده برای مقاوم‌سازی ویسپر-کوچک مؤثرتر است یا افزودن بهسازی و همجوشی دوجریانه. این مسئله با تخریب کنترل‌شده و مقایسهٔ آموزش چندشرایطی با همجوشی آغاز و با آزمون افزایش گفتار واقعی از صوت‌های بلند ایران‌صدا دنبال می‌شود.

\قسمت{اهمیت و دامنهٔ پژوهش}

تماس‌های تلفنی منبع مهمی از داده‌های گفتاری‌اند، اما کانال مخابراتی اطلاعات آکوستیکی را تغییر می‌دهد. رونویسی تماس‌های واقعی فارسی نیز پرهزینه و همراه با ملاحظات حقوقی و حریم خصوصی است. شبیه‌سازی کانال می‌تواند این کمبود را کاهش دهد، به‌شرط تولید بازتولیدپذیر و ارزیابی روی دادهٔ واقعی. گفتار بلند وب نیز تنوع را افزایش می‌دهد، اما به قطعه‌بندی، کنترل کیفیت، حفظ منشأ و جلوگیری از نشت نیاز دارد.

برای مرحلهٔ نهایی، حدود ۳۷۰۰ کتاب صوتی، معادل نزدیک به ۷۰۰۰ ساعت صوت خام، از ایران‌صدا گردآوری شده و حدود ۱۰۰۰ ساعت برای آزمایش انتخاب می‌شود. قطعه‌های ۱۵ تا ۲۵ ثانیه‌ای با آشکارسازی فعالیت گفتاری سیلرو ساخته و با ویسپر-متوسطِ سازگارشده با فارسی شبه‌برچسب‌گذاری می‌شوند. پس از نرمال‌سازی قطعی، اصلاح نگارشی محدود با مدل زبانی بزرگ انجام خواهد شد، بدون افزودن، حذف، استنباط، جابه‌جایی یا بازنویسی محتوا. متن‌ها، زمان‌ها، شناسهٔ منبع، اثر انگشت و تنظیمات نگه‌داری می‌شوند تا ممیزی ممکن باشد.

تمرکز پژوهش بر دقت بازشناسی است، نه کیفیت ادراکی صدا. بهساز باید اطلاعات بازشناسی را حفظ کند؛ زیرا ممکن است با وجود بهبود شنیداری، مصنوعات تازه ایجاد کند \مرجع{iwamoto2022bad,iwamoto2024artifacts}. ازاین‌رو، مدل هر دو نمای اصلی و بهسازی‌شده را دریافت می‌کند. آزمایش‌هایی با فست کانفورمر\پاورقی{FastConformer} نیز انجام شد، اما سامانهٔ نهایی بر ویسپر-کوچک تمرکز دارد.

\قسمت{اهداف پژوهش}

هدف اصلی، طراحی و ارزیابی مسیری بازتولیدپذیر و داده‌محور برای افزایش مقاومت بازشناسی گفتار فارسی در دامنه‌های عمومی، تلفنی و VoIP است. اهداف جزئی عبارت‌اند از:

\شروع{فقرات}
\فقره یکپارچه‌سازی منابع موجود گفتار فارسی، ساخت نمونه‌های بلندتر و تولید جفت‌های هم‌تراز پاک--تخریب‌شده با ثبت بذر و فرادادهٔ نویز، بهره، پهنای‌باند، چرخهٔ واقعی کدک و افت بسته؛
\فقره سنجش اثر آموزش چندشرایطی با مقایسهٔ کنترل‌شدهٔ ویسپر-کوچک آموزش‌دیده با داده‌های معمولی و همان مدل آموزش‌دیده با داده‌های معمولی، بلند و تخریب‌شده؛
\فقره طراحی و ارزیابی مدل دوجریانه‌ای که نمای اصلی و بهسازی‌شدهٔ لاگ-مل را با رمزگذار مشترک ویسپر، توجه متقاطع دوسویه و دروازهٔ آموختنی ترکیب می‌کند؛
\فقره ساخت خط لوله‌ای ردیابی‌پذیر برای تبدیل صوت بلند ایران‌صدا به قطعه‌های پالایش‌شده و شبه‌برچسب‌خورده، با تقسیم در سطح کتاب صوتی؛
\فقره مقایسهٔ نهایی خط مبنای چندشرایطی با همان معماری و مخلوط کامل آموزشی به‌علاوهٔ دادهٔ پذیرفته‌شدهٔ ایران‌صدا، روی مجموعه‌های آزمون ثابت و انسانی؛ و
\فقره گزارش نرخ خطای واژه و نویسه به تفکیک گفتار عمومی و واقعی تلفنی/VoIP و تحلیل صریح نتایج مثبت، خنثی یا منفی و محدودیت‌های آن‌ها.
\پایان{فقرات}

\قسمت{پرسش‌های پژوهش}

آزمایش‌های این پژوهش برای پاسخ به شش پرسش زیر طراحی شده‌اند:

\شروع{فقرات}
\فقره آیا افزودن گفتار مصنوعی تخریب‌شده به آموزش، خطای ویسپر-کوچک را نسبت به ریزتنظیم روی داده‌های معمولی کاهش می‌دهد؟
\فقره آیا مزیت آموزش چندشرایطی به گفتار واقعی تلفنی و VoIP منتقل می‌شود، بی‌آن‌که عملکرد روی گفتار عمومی فارسی به‌طور چشمگیری افت کند؟
\فقره آیا همجوشی نمای اصلی و بهسازی‌شده می‌تواند از ویسپر معمولی و، مهم‌تر از آن، خط مبنای چندشرایطی بهتر عمل کند؟
\فقره آیا دروازهٔ آموختنی از هر دو نما به‌طور معنادار استفاده می‌کند یا عملاً به یک جریان متکی می‌شود؟
\فقره آیا افزودن گفتار شبه‌برچسب‌خوردهٔ منتخب ایران‌صدا به مخلوط موجودِ داده‌های معمولی، بلند و تخریب‌شده، عملکرد ویسپر-کوچک را روی همان آزمون‌های انسانی بهبود می‌دهد؟
\فقره سود یا زیان حاصل از دادهٔ ایران‌صدا میان گفتار عمومی فارسی و گفتار واقعی تلفنی/VoIP چگونه تغییر می‌کند؟
\پایان{فقرات}

چهار پرسش نخست با آزمایش‌های آموزش چندشرایطی و همجوشی بررسی شده‌اند و پرسش‌های پنجم و ششم پس از تکمیل پالایش داده، آموزش و ارزیابی نهایی پاسخ داده می‌شوند. پاسخ همهٔ پرسش‌ها در فصل نتایج و بر پایهٔ نرخ خطای واژه و نویسه، عملکرد روی هر مجموعه‌داده و تحلیل رفتار دروازه گزارش خواهد شد؛ در این فصل دربارهٔ جهت یا اندازهٔ نتایج پیش‌داوری نمی‌شود.

\قسمت{دستاوردها و حدود ادعا}

دستاوردهای پژوهش یک مسیر آزمایشی پیوسته می‌سازند: آماده‌سازی نزدیک به هزار ساعت گفتار موجود در قالب مشترک؛ ساخت نمونه‌های بلند و خط لولهٔ بازتولیدپذیر تخریب؛ مقایسهٔ کنترل‌شدهٔ آموزش معمولی و چندشرایطی؛ و طراحی و ارزیابی معماری دوجریانه. این طراحی امکان می‌دهد اثر افزایش تنوع داده از اثر افزودن پیچیدگی معماری جدا شود و رفتار همجوشی نیز در برابر یک خط مبنای چندشرایطی سنجیده شود.

مرحلهٔ نهایی، یک خط لولهٔ بازتولیدپذیر برای گردآوری، قطعه‌بندی، شبه‌برچسب‌گذاری، نرمال‌سازی، اصلاح محافظه‌کارانه و پالایش صوت بلند ایران‌صدا و یک آزمایش کنترل‌شدهٔ افزایش مقیاس داده است. چون اصلاح زبانی، آمار نهایی پیکره، آموزش ویسپر-کوچک و ارزیابی آن هنوز تکمیل نشده‌اند، تنها به‌عنوان روش و آزمایش برنامه‌ریزی‌شده گزارش می‌شوند. ادعاهای نهایی به مقایسه‌های کنترل‌شده روی مجموعه‌های آزمون ثابت محدود خواهند بود و به همهٔ کانال‌ها، گویش‌ها و منابع فارسی تعمیم داده نمی‌شوند.

\قسمت{ساختار پایان‌نامه}

این پایان‌نامه در شش فصل تنظیم شده است. فصل دوم، مفاهیم بازشناسی گفتار فارسی، ویسپر، ناهمخوانی دامنه، آموزش چندشرایطی، بهسازی، گفتار وب، قطعه‌بندی صوت بلند، شبه‌برچسب‌گذاری و اصلاح محافظه‌کارانهٔ متن را توضیح می‌دهد. فصل سوم پژوهش‌های پیشین در بازشناسی فارسی، مقاومت در برابر تخریب، خودآموزی، ساخت پیکرهٔ وب و همجوشی را مرور می‌کند. فصل چهارم داده‌های اولیه، خط لولهٔ تخریب، ریزتنظیم ویسپر، معماری دوجریانه و فرایند ایران‌صدا از گردآوری تا آزمایش نهایی را شرح می‌دهد. فصل پنجم نتایج تکمیل‌شدهٔ تخریب و همجوشی و، پس از آماده‌شدن، آمار پیکره و مقایسهٔ ثابت مدل ایران‌صدا را همراه با تهدیدهای اعتبار گزارش می‌کند. سرانجام، فصل ششم به جمع‌بندی شواهد نهایی، محدودیت‌ها و مسیرهای آینده اختصاص دارد.
